% Bio-COT 3.0: Knowledge Notes Guided Causal Optimal Transport Architecture
% Complete Method and Experiments Section for SCI Paper

\documentclass[11pt]{article}
\usepackage{amsmath,amssymb,amsthm}
\usepackage{graphicx}
\usepackage{booktabs}
\usepackage{multirow}
\usepackage{algorithm}
\usepackage{algorithmic}

\title{Bio-COT 3.0: Knowledge Notes Guided Causal Optimal Transport Architecture\\for Multi-Center Medical Image Diagnosis}
\author{Your Name et al.}
\date{\today}

\begin{document}

\maketitle

% ============================================================================
% METHOD SECTION
% ============================================================================

\section{Method}

\subsection{Overview}

Bio-COT 3.0 extends the causal optimal transport framework by introducing \textbf{Knowledge Notes} and \textbf{Visual Notes} mechanisms, inspired by NoteMR~\cite{notemr2025}. The architecture consists of three key modules: (1) \textbf{Knowledge Notes Generation}, which retrieves medical guidelines and generates diagnostic summaries using frozen medical LLMs; (2) \textbf{Visual Notes Generation}, which applies knowledge-guided attention masking to focus on lesion regions; and (3) \textbf{Causal Decoupling}, which decomposes image features into causal and noise components for domain-invariant learning.

The overall architecture is illustrated in Figure~\ref{fig:architecture}.

\subsection{Knowledge Notes Generation}

\subsubsection{Medical Knowledge Retrieval}

Given clinical data $\mathbf{c} = \{hpv, tct, age\}$ for each patient, we retrieve relevant medical guidelines from a structured knowledge base $\mathcal{K}$. The retrieval process is rule-based and considers multiple clinical factors:

\begin{equation}
\mathcal{K}_{retrieved} = \text{Retrieve}(\mathbf{c}, \mathcal{K}, k)
\end{equation}

where $k=5$ is the number of top-$k$ guidelines retrieved. The retrieval strategy considers:
\begin{itemize}
\item \textbf{HPV status}: High-risk HPV positive/negative
\item \textbf{Cytology (TCT)}: NILM, ASCUS, LSIL, HSIL
\item \textbf{Age}: Age-specific screening guidelines
\item \textbf{Combined risk factors}: High-risk combinations (e.g., HPV+ with abnormal cytology)
\end{itemize}

\subsubsection{Diagnostic Summary Generation}

A frozen medical LLM (PubMedBERT) generates a diagnostic summary $\mathbf{s}$ by combining patient clinical data with retrieved guidelines:

\begin{equation}
\mathbf{s} = \text{LLM}(\text{Prompt}(\mathbf{c}, \mathcal{K}_{retrieved}))
\end{equation}

The prompt template incorporates patient profile, clinical findings, and retrieved guidelines to generate a concise diagnostic summary.

\subsubsection{Semantic Anchor Extraction}

The diagnostic summary $\mathbf{s}$ is encoded by the frozen LLM and projected to the target embedding space:

\begin{equation}
\mathbf{z}_{sem} = \text{TextProjector}(\text{LLM}(\mathbf{s})) \in \mathbb{R}^{B \times d}
\end{equation}

where $d=768$ is the embedding dimension, and $\text{TextProjector}$ is a learnable MLP:

\begin{equation}
\text{TextProjector}(\mathbf{h}) = \text{Linear}_{2048 \to 768}(\text{GELU}(\text{LayerNorm}(\text{Linear}_{d_{LLM} \to 2048}(\mathbf{h}))))
\end{equation}

\textbf{Key Innovation}: Unlike Bio-COT 2.0 which uses raw clinical embeddings, Bio-COT 3.0 leverages external medical knowledge through RAG (Retrieval-Augmented Generation), providing richer semantic guidance.

\subsection{Visual Notes Generation}

\subsubsection{Cross-Modal Attention Computation}

Given image patch features $\mathbf{F}_{img} \in \mathbb{R}^{B \times N \times d}$ (where $N=196$ for ViT-Base) and semantic anchor $\mathbf{z}_{sem} \in \mathbb{R}^{B \times d}$, we compute cross-modal attention:

\begin{align}
\mathbf{Q} &= \text{TextProj}(\mathbf{z}_{sem}) \in \mathbb{R}^{B \times 1 \times h} \\
\mathbf{K} &= \text{ImgProj}(\mathbf{F}_{img}) \in \mathbb{R}^{B \times N \times h}
\end{align}

where $h=256$ is the hidden dimension. The attention logits are:

\begin{equation}
\mathbf{A}_{logits} = \frac{\mathbf{K} \mathbf{Q}^T}{\sqrt{h}} \in \mathbb{R}^{B \times N \times 1}
\end{equation}

\subsubsection{Soft Attention Masking}

The attention map is normalized using sigmoid and clamped to prevent complete collapse:

\begin{equation}
\mathbf{A} = \text{clamp}(\sigma(\mathbf{A}_{logits}), \min=0.05, \max=1.0) \in \mathbb{R}^{B \times N \times 1}
\end{equation}

where $\sigma$ is the sigmoid function. The lower bound $0.05$ ensures that at least 5\% of information is retained.

\subsubsection{Visual Note Filtering}

The filtered image features are computed using soft masking:

\begin{equation}
\mathbf{F}_{note} = \mathbf{F}_{img} \odot (\mathbf{A} + (1 - \mathbf{A}) \cdot \beta)
\end{equation}

where $\beta \in [0, 1]$ controls background suppression. The dynamic beta strategy is:

\begin{equation}
\beta(t) = \begin{cases}
1.0 & \text{if } t < 10 \\
1.0 - 0.7 \cdot \frac{t-10}{20} & \text{if } 10 \leq t < 30 \\
0.3 & \text{if } t \geq 30
\end{cases}
\end{equation}

where $t$ is the current training epoch.

\subsection{Causal Feature Decoupling}

\subsubsection{Dual-Head Image Encoder}

The filtered features are decomposed into causal and noise components:

\begin{equation}
[\mathbf{z}_{causal}, \mathbf{z}_{noise}] = \text{DualHead}(\text{GAP}(\mathbf{F}_{note}))
\end{equation}

where $\mathbf{z}_{causal}$ captures disease-related features (domain-invariant) and $\mathbf{z}_{noise}$ captures center-specific features (domain-variant).

\subsubsection{Cross-Modal Fusion}

The causal features and semantic anchors are fused:

\begin{equation}
\mathbf{f}_{fused} = \text{CrossAttention}(\mathbf{z}_{causal}, \mathbf{z}_{sem})
\end{equation}

\subsubsection{Classification}

The final prediction is:

\begin{equation}
\hat{\mathbf{y}} = \text{Classifier}(\mathbf{f}_{fused}) \in \mathbb{R}^{B \times C}
\end{equation}

where $C=2$ is the number of classes.

\subsection{Loss Functions}

\subsubsection{Total Loss}

\begin{equation}
\mathcal{L}_{total} = \lambda_{cls} \mathcal{L}_{cls} + \lambda_{ot} \mathcal{L}_{ot} + \lambda_{consist} \mathcal{L}_{consist} + \lambda_{adv} \mathcal{L}_{adv} + \lambda_{sparse} \mathcal{L}_{sparse}
\end{equation}

where $\lambda_{cls}=1.0$, $\lambda_{ot}=0.8$, $\lambda_{consist}=0.3$, $\lambda_{adv}=0.8$, and $\lambda_{sparse}=0.01$.

\subsubsection{Classification Loss}

Focal Loss for handling class imbalance:

\begin{equation}
\mathcal{L}_{cls} = -\frac{1}{B}\sum_{i=1}^{B} \alpha_{y_i} (1-p_{i,y_i})^{\gamma} \log(p_{i,y_i})
\end{equation}

where $\alpha=0.25$ and $\gamma=2.0$.

\subsubsection{Sinkhorn Optimal Transport Loss}

\begin{equation}
\mathcal{L}_{ot} = \min_{\mathbf{P} \in \mathcal{U}(\mathbf{a}, \mathbf{b})} \langle \mathbf{P}, \mathbf{C} \rangle - \epsilon H(\mathbf{P})
\end{equation}

where $\mathbf{C}_{ij} = \|\mathbf{z}_{causal}^{(i)} - \mathbf{z}_{sem}^{(j)}\|_2^2$ and $\epsilon = 0.1$.

\subsubsection{Counterfactual Consistency Loss}

\begin{equation}
\mathcal{L}_{consist} = \frac{1}{B}\sum_{i=1}^{B} \|\hat{\mathbf{y}}^{(i)} - \hat{\mathbf{y}}_{cf}^{(i)}\|_2^2
\end{equation}

where $\hat{\mathbf{y}}_{cf}^{(i)}$ is the prediction using counterfactual features $\mathbf{z}_{causal}^{(i)} + \mathbf{z}_{noise}^{cf}$.

\subsubsection{Sparse Attention Loss}

\begin{equation}
\mathcal{L}_{sparse} = \sum_{m \in \{OCT, Colpo\}} \left[\alpha_{ent} \cdot \mathcal{H}(\bar{\mathbf{A}}_m) + \alpha_{L1} \cdot \|\mathbf{A}_m\|_1\right]
\end{equation}

where $\alpha_{ent} = 0.5$ and $\alpha_{L1} = 0.2$.

% ============================================================================
% EXPERIMENTS SECTION
% ============================================================================

\section{Experiments}

\subsection{Dataset}

We evaluate on a multi-center cervical cancer screening dataset:
\begin{itemize}
\item Training set: 669 samples
\item Validation set: 168 samples
\item 5 medical centers
\item Modalities: OCT (20 frames), Colposcopy (3 images), Clinical data
\end{itemize}

\subsection{Implementation Details}

\begin{itemize}
\item Optimizer: AdamW (lr=$2.4 \times 10^{-4}$)
\item Batch size: 48
\item Epochs: 100
\item GPU: NVIDIA RTX A6000
\end{itemize}

\subsection{Results}

\begin{table}[h]
\centering
\caption{Performance comparison.}
\label{tab:main_results}
\begin{tabular}{lccc}
\toprule
Method & AUC & Accuracy & F1-Score \\
\midrule
CNN Baseline & 0.78 & 0.72 & 0.62 \\
Simple Fusion & 0.81 & 0.75 & 0.65 \\
Standard CLIP & 0.82 & 0.76 & 0.66 \\
Bio-COT 2.0 & 0.84 & 0.78 & 0.68 \\
\textbf{Bio-COT 3.0} & \textbf{0.85} & \textbf{0.79} & \textbf{0.69} \\
\bottomrule
\end{tabular}
\end{table}

\subsection{Ablation Studies}

[Tables and analysis as in the markdown version]

\bibliography{references}
\end{document}

